\chapter{随机向量的数字特征}
\section{随机变量的期望与方差}
\begin{definition}[离散型随机变量的数学期望]
设随机变量 \(X\) 是离散型随机变量，它的概率函数为
\[
P(X=x_k)=p_k, \quad k=1,2,\ldots
\]
如果
\[
\sum_{k=1}^{\infty} |x_k|p_k < \infty,
\]
则称 \(X\) 的\textbf{数学期望}为
\[
E(X)=\sum_{k=1}^{\infty} x_k p_k.
\]
\end{definition}
\begin{example}
开学伊始，某宿舍楼有 $n$ 个房间，每个房间只有一把没有标记的钥匙放在新来的管理员手中。  
某同学不慎丢掉钥匙，向管理员求助，求管理员开门所需次数的数学期望。

\begin{solution}
设试开次数为 $X$，则  
\[
P(X=k)=\frac{1}{n}, \quad k=1,2,\ldots,n.
\]
于是
\[
E(X)
=\sum_{k=1}^{n} k \cdot \frac{1}{n}
=\frac{1}{n}\cdot \frac{(1+n)n}{2}
=\frac{n+1}{2}.
\]
\end{solution}
\end{example}
\begin{note}
不难计算得：

\begin{itemize}
  \item 若 \(X\) 服从参数为 \(\lambda\) 的泊松分布，则
  \[
  E(X)=\lambda.
  \]

  \item 若 \(X\) 服从参数为 \(n,p\) 的二项分布，则
  \[
  E(X)=np.
  \]
\end{itemize}
\end{note}
\begin{definition}[连续型随机变量的数学期望---概率f（x）加权重x的绝对值]
设随机变量 \(X\) 是连续型随机变量，其概率密度函数为 \(f(x)\)。  
如果
\[
\int_{-\infty}^{\infty} |x| f(x)\,dx < \infty,
\]
则称 \(X\) 的\textbf{数学期望}为
\[
E(X)=\int_{-\infty}^{\infty} x f(x)\,dx.
\]

也就是说，连续型随机变量的数学期望是一个绝对可积的积分。
\end{definition}
\begin{note}
不难计算得：

\begin{itemize}
  \item 若 \(X \sim U(a,b)\)，即 \(X\) 服从 \((a,b)\) 上的均匀分布，则
  \[
  E(X)=\frac{a+b}{2}.
  \]

  \item 若 \(X \sim N(\mu,\sigma^2)\)，则
  \[
  E(X)=\mu.
  \]
\end{itemize}
\end{note}
\begin{example}
已知某地区成年男子身高 \(X \sim N(1.68, \sigma^2)\)，则
\[
E(X)=\mu=1.68.
\]

这意味着，若从该地区抽查很多个成年男子，分别测量他们的身高，那么这些身高的平均值近似是 \(1.68\)。
\end{example}
\begin{definition}[随机变量函数的数学期望]
设 \(X\) 是一个随机变量，\(Y=g(X)\)。若

\begin{itemize}
  \item 当 \(X\) 为离散型时，\(P(X=x_k)=p_k\)；
  \item 当 \(X\) 为连续型时，\(X\) 的概率密度函数为 \(f(x)\)，
\end{itemize}

则随机变量 \(Y=g(X)\) 的数学期望为
\[
E(Y)=E[g(X)]=
\begin{cases}
  \displaystyle \sum_{k=1}^{\infty} g(x_k)p_k, & X\text{ 为离散型},\\[1em]
  \displaystyle \int_{-\infty}^{\infty} g(x)f(x)\,dx, & X\text{ 为连续型}.
\end{cases}
\]
\end{definition}
\begin{theorem}[数学期望的性质]
设 \(X, Y, X_1, X_2, \ldots, X_n\) 为随机变量，\(C, k\) 为常数，则有以下性质：

\begin{enumerate}
  \item 设 \(C\) 是常数，则 \(E(C)=C\)；
  
  \item 若 \(k\) 是常数，则 \(E(kX)=kE(X)\)；
  
  \item \(E(X_1+X_2)=E(X_1)+E(X_2)\)，  
  推广为：
  \[
  E\!\left[\sum_{i=1}^{n} X_i\right]=\sum_{i=1}^{n}E(X_i)；
  \]
  
  \item 若 \(X, Y\) 独立，则 \(E(XY)=E(X)E(Y)\)，  
  推广为：
  \[
  E\!\left[\prod_{i=1}^{n} X_i\right]
  =\prod_{i=1}^{n}E(X_i)\quad (\text{当诸 }X_i\text{ 独立时})。
  \]
\end{enumerate}

\begin{remark}
注意：由 \(E(XY)=E(X)E(Y)\) \textbf{不一定能推出} \(X,Y\) 独立。
\end{remark}
\end{theorem}
\begin{definition}[二维随机变量的期望定义]
对于二维连续型随机变量 $(X,Y)$，期望定义推广为：
\[
E[g(X,Y)] = \iint_{\mathbb{R}^2} g(x,y)\, f(x,y)\, dxdy,
\]
其中 $g(x,y)$ 是 $X,Y$ 的某个函数。
\end{definition}

\begin{note}[我们关心的三种特殊情形]
\begin{enumerate}
  \item 对于 $E(X)$，取 $g(x,y)=x$，因此
  \[
  E(X) = \iint x f(x,y)\,dxdy.
  \]

  \item 对于 $E(Y)$，取 $g(x,y)=y$，因此
  \[
  E(Y) = \iint y f(x,y)\,dxdy.
  \]

  \item 对于 $E(XY)$，取 $g(x,y)=xy$，因此
  \[
  E(XY) = \iint xy f(x,y)\,dxdy.
  \]
\end{enumerate}
\end{note}
\begin{note}[积分中的“平均”思想]
关键在于：
\[
\iint_{\mathbb{R}^2} f(x,y)\,dxdy = 1.
\]

这意味着——  
整个积分的“权重”总和是 \(1\)。

所以这个积分相当于：
\[
E[g(X,Y)] 
= \frac{\displaystyle \iint g(x,y)\,f(x,y)\,dxdy}
{\displaystyle \iint f(x,y)\,dxdy}.
\]

也就是说，期望其实就是在概率密度 \(f(x,y)\) 下，
对函数 \(g(x,y)\) 做的加权平均。
\end{note}
\section{方差}
\begin{definition}[方差与标准差]
设 $X$ 是一个随机变量，若 
\[
E[(X - E(X))^2] < \infty,
\]
则称
\[
D(X) = E\!\left[(X - E(X))^2\right]
\]
为随机变量 $X$ 的\textbf{方差}。

方差的算术平方根
\[
\sqrt{D(X)}
\]
称为随机变量 $X$ 的\textbf{标准差}。
\end{definition}
这章重点是计算级数的计算

\begin{note}方差的意义:
方差刻画了随机变量取值的\textbf{离散程度}。

\[
D(X) = E\!\left[(X - E(X))^2\right].
\]

若随机变量 $X$ 的取值比较集中，则方差较小；  
若随机变量 $X$ 的取值比较分散，则方差较大。

特别地，当方差 $D(X)=0$ 时，随机变量 $X$ 以概率 $1$ 取常数值。
\end{note}
\begin{theorem}[方差的简化公式]
\[
D(X) = E(X^2) - [E(X)]^2.
\]
\end{theorem}

\begin{proof}
由方差定义：
\[
D(X) = E\!\left[(X - E(X))^2\right].
\]

展开平方项：
\[
D(X) = E\!\left[X^2 - 2X E(X) + [E(X)]^2\right].
\]

利用期望的线性性质：
\[
D(X) = E(X^2) - 2E(X)E(X) + [E(X)]^2.
\]

整理得：
\[
\boxed{D(X) = E(X^2) - [E(X)]^2.}
\]
\end{proof}
\begin{note}
    E（x）是一个数可以提出来
\end{note}
\begin{theorem}[方差的性质]
设随机变量 $X, X_1, X_2, \ldots, X_n$，常数 $C, C_i$，则有以下性质：

\begin{enumerate}
  \item 若 $C$ 为常数，则
  \[
  D(C)=0.
  \]
  
  \item 若 $C$ 为常数，则
  \[
  D(CX)=C^2 D(X).
  \]
  
  \item 若 $X_1$ 与 $X_2$ 相互独立，则
  \[
  D(X_1+X_2)=D(X_1)+D(X_2).
  \]
  
  更一般地，若 $X_1, X_2, \ldots, X_n$ 相互独立，则
  \[
  D\!\left(\sum_{i=1}^n X_i\right)
  = \sum_{i=1}^n D(X_i),
  \qquad
  D\!\left(\sum_{i=1}^n C_i X_i\right)
  = \sum_{i=1}^n C_i^2 D(X_i).
  \]
\end{enumerate}
\end{theorem}

\begin{note}
若 $X_1$ 与 $X_2$ 不一定独立，则
\[
D(X_1+X_2) = D(X_1) + D(X_2) + 2\,\mathrm{Cov}(X_1,X_2),
\]
其中 $\mathrm{Cov}(X_1,X_2)=E[(X_1-E(X_1))(X_2-E(X_2))]$ 为协方差。
\end{note}


\begin{example}[二项分布的方差]
设 $X \sim B(n,p)$，则 $X$ 表示 $n$ 重贝努里试验中的“成功”次数。

若设
\[
X_i=
\begin{cases}
1, & \text{如第 $i$ 次试验成功},\\[4pt]
0, & \text{如第 $i$ 次试验失败},
\end{cases}
\qquad i=1,2,\ldots,n,
\]
则有
\[
X=\sum_{i=1}^{n} X_i
\]
表示 $n$ 次试验中的“成功”次数。

由于 $E(X_i)=P(X_i=1)=p$，且
\[
E(X_i^2)=1^2\cdot p+0^2\cdot(1-p)=p,
\]
所以
\[
D(X_i)=E(X_i^2)-[E(X_i)]^2=p-p^2=p(1-p).
\]

因为 $X_1,X_2,\ldots,X_n$ 相互独立，
\[
D(X)=D\!\left(\sum_{i=1}^{n}X_i\right)=\sum_{i=1}^{n}D(X_i)=np(1-p).
\]
\end{example}
\begin{theorem}[方差的性质（标准化与最小化）]
设随机变量 $X$ 的方差存在。

\begin{enumerate}[(1)]
  \item 若令
  \[
  Y=\frac{X-E(X)}{\sqrt{D(X)}},
  \]
  则称 $Y$ 为 $X$ 的标准化变量，且有
  \[
  E(Y)=0,\qquad D(Y)=1.
  \]

  \item 对任意实数 $c$，随机变量 $X-c$ 的方差亦存在，且成立不等式
  \[
  D(X)\le E\!\big([X-c]^2\big),
  \]
  等号当且仅当 $c=E(X)$ 时取到。
\end{enumerate}
\end{theorem}
\begin{definition}[切比雪夫不等式]
设随机变量 $X$ 有期望 $E(X)$ 和方差 $\sigma^2$，  
则对任意 $\varepsilon>0$，有
\[
P\big(|X-E(X)|\ge\varepsilon\big)\le\frac{\sigma^2}{\varepsilon^2},
\]
或等价地，
\[
P\big(|X-E(X)|<\varepsilon\big)\ge 1-\frac{\sigma^2}{\varepsilon^2}.
\]

\end{definition}

\begin{note}
由切比雪夫不等式可知：若方差 $\sigma^2$ 越小，  
则事件 $\{|X-E(X)|<\varepsilon\}$ 的概率越大，  
即随机变量 $X$ 的取值越集中在其期望附近。
\end{note}
\section{矩}

\begin{definition}
设 $X$ 是随机变量，$k$ 是非负整数。
\begin{itemize}
  \item 如果 $E(X^k)$ 存在，记作 $\nu_k$，称为 $X$ 的 $k$ 阶原点矩；
  \item 如果 $E(|X|^k)$ 存在，记作 $\alpha_k$，称为 $X$ 的 $k$ 阶原点绝对矩；
  \item 如果 $E\{[X-E(X)]^k\}$ 存在，记作 $\mu_k$，称为 $X$ 的 $k$ 阶中心矩；
  \item 如果 $E|X-E(X)|^k$ 存在，记作 $\beta_k$，称为 $X$ 的 $k$ 阶中心绝对矩。
\end{itemize}
\end{definition}
\begin{theorem}[Markov 不等式]
设 $X$ 是随机变量，$k$ 为自然数。若 $E(X^k)$ 存在，则对任意 $\varepsilon > 0$，有
\[
P(|X|>\varepsilon)\le \frac{E(|X|^k)}{\varepsilon^k}.
\]
\end{theorem}
\begin{theorem}
设 $X, Y$ 是随机变量。若 $E(X^2)$ 与 $E(Y^2)$ 皆存在，  
则 $E(XY)$ 亦存在，且有
\[
\bigl[E(|XY|)\bigr]^2 \le E(X^2)\,E(Y^2).
\]
\end{theorem}
\begin{theorem}
设 $X$ 是随机变量，$n$ 为自然数。若 $E(X^n)$ 存在，  
则对任意正整数 $k \le n$，$E(X^k)$ 亦存在。
\end{theorem}
\section{协方差}
\begin{definition}
任意两个随机变量 $X$ 和 $Y$ 的协方差，记为 $\mathrm{Cov}(X,Y)$，定义为
\[
\mathrm{Cov}(X,Y) = E\{[X - E(X)][Y - E(Y)]\}.
\]
\end{definition}
\begin{theorem}[协方差的简单性质如下]：
\begin{itemize}
  \item[(1)] $\mathrm{Cov}(X,Y) = \mathrm{Cov}(Y,X)$；
  \item[(2)] 若 $a,b$ 是常数，则 $\mathrm{Cov}(aX,bY) = ab\,\mathrm{Cov}(X,Y)$；
  \item[(3)] $\mathrm{Cov}(X_1 + X_2, Y) = \mathrm{Cov}(X_1, Y) + \mathrm{Cov}(X_2, Y)$。
\end{itemize}
\end{theorem}
\begin{proposition}[计算协方差的一个简单公式]
由协方差的定义及期望的性质，可得
\[
\begin{aligned}
\mathrm{Cov}(X,Y)
&=E\{[X-E(X)][Y-E(Y)]\} \\[4pt]
&=E(XY)-E(X)E(Y)-E(Y)E(X)+E(X)E(Y) \\[4pt]
&=E(XY)-E(X)E(Y).
\end{aligned}
\]
因此，
\[
\mathrm{Cov}(X,Y)=E(XY)-E(X)E(Y).
\]
\end{proposition}
\begin{note}
可见，若 $X$ 与 $Y$ 独立，则 $\mathrm{Cov}(X,Y)=0$。

协方差的大小在一定程度上反映了 $X$ 和 $Y$ 相互间的关系，但由于
\[
\mathrm{Cov}(kX,kY)=k^2\mathrm{Cov}(X,Y),
\]
协方差的数值会随度量单位的改变而改变，因此仅用协方差来衡量变量间的相关程度并不合理。
\end{note}

\begin{definition}
设 $D(X) > 0$，$D(Y) > 0$，称
\[
\rho_{XY} = \frac{\mathrm{Cov}(X,Y)}{\sqrt{D(X)D(Y)}}
\]
为随机变量 $X$ 和 $Y$ 的相关系数。

在不致引起混淆时，记 $\rho_{XY}$ 为 $\rho$。
\end{definition}
\begin{proposition}[相关系数的性质]
\begin{itemize}
  \item[(1)] 若 $\rho$ 为随机变量 $X$ 与 $Y$ 的相关系数，则有
  \[
  |\rho| \le 1.
  \]
  证明：由方差的性质和协方差的定义知，对任意实数 $x$，有
  \[
  0 \le D(Y - xX) = x^2 D(X) - 2x\,\mathrm{Cov}(X,Y) + D(Y),
  \]
  由此可得 $|\rho| \le 1$。

  \item[(2)] 当 $X$ 与 $Y$ 独立时，$\rho = 0$，但其逆不真。  
  由于 $X$ 与 $Y$ 独立时，$\mathrm{Cov}(X,Y)=0$，故
  \[
  \rho = \frac{\mathrm{Cov}(X,Y)}{\sqrt{D(X)D(Y)}} = 0.
  \]
  但由 $\rho = 0$ 并不能推出 $X$ 与 $Y$ 独立。
\end{itemize}
\end{proposition}
\begin{note}
    在二维正态分布中不相关等于独立，其他的不等价独立可推不相关不相关不能推正态分布。
\end{note}
\begin{example}[相关系数为 $0$ 但不独立的例子]
设随机变量 $X$ 在区间 $[-1,1]$ 上服从均匀分布，即
\[
f_X(x)=
\begin{cases}
\dfrac{1}{2}, & -1\le x\le 1,\\[3pt]
0, & \text{其他}.
\end{cases}
\]
定义 $Y = X^2$。

显然，$Y$ 完全由 $X$ 决定，因此 $X$ 与 $Y$ 不独立。

计算它们的相关系数：
\[
E(X)=0,\qquad E(Y)=E(X^2)=\int_{-1}^1 x^2 \cdot \frac{1}{2}\,dx = \frac{1}{3}.
\]
又有
\[
E(XY)=E(X^3)=\int_{-1}^1 x^3 \cdot \frac{1}{2}\,dx = 0.
\]
于是协方差为
\[
\mathrm{Cov}(X,Y)=E(XY)-E(X)E(Y)=0-0\times\frac{1}{3}=0.
\]
因此，
\[
\rho=\frac{\mathrm{Cov}(X,Y)}{\sqrt{D(X)D(Y)}}=0.
\]
可见，虽然 $\rho=0$，但由于 $Y=X^2$ 与 $X$ 之间存在确定的函数关系，$X$ 与 $Y$ 并不独立。
\end{example}
\begin{theorem}
设 $X,Y$ 为随机变量，且 $E(X^2),E(Y^2)$ 存在。则下述四个命题等价：
\begin{enumerate}[(1)]
  \item $\mathrm{Cov}(X,Y)=0$；
  \item $\rho_{XY}=0$（$X$ 与 $Y$ 不相关）；
  \item $E(XY)=E(X)\,E(Y)$；
  \item $D(X+Y)=D(X)+D(Y)$。
\end{enumerate}
\end{theorem}
\begin{proposition}
若相关系数满足 $|\rho| = 1$，则当且仅当存在常数 $a,b$（其中 $b \ne 0$），使得
\[
P\{Y = a + bX\} = 1。
\]
即 $X$ 和 $Y$ 以概率 $1$ 线性相关。

因此，相关系数刻画了随机变量 $X$ 与 $Y$ 之间“线性相关”的程度。
\end{proposition}
\begin{definition}
设 $(X_1, X_2)$ 为二维随机变量，定义其四个二阶中心矩如下：
\[
\begin{aligned}
c_{11} &= E\{[X_1 - E(X_1)]^2\}, \\[3pt]
c_{12} &= E\{[X_1 - E(X_1)][X_2 - E(X_2)]\}, \\[3pt]
c_{21} &= E\{[X_2 - E(X_2)][X_1 - E(X_1)]\}, \\[3pt]
c_{22} &= E\{[X_2 - E(X_2)]^2\}.
\end{aligned}
\]

将其排列成矩阵形式：
\[
C =
\begin{pmatrix}
c_{11} & c_{12} \\[3pt]
c_{21} & c_{22}
\end{pmatrix},
\]
称矩阵 $C$ 为随机向量 $(X_1, X_2)$ 的\textbf{协方差矩阵}。

由于 $c_{12} = c_{21}$，该矩阵是一个对称矩阵。
\end{definition}

% 条件数学期望的数学期望（重期望）

若把条件数学期望 $g(y) = E(X \mid Y = y)$ 中的 $y$ 换成随机变量 $Y$，
并记为 $g(Y) = E(X \mid Y)$，则其为随机变量，其数学期望称为 **重期望**，
即
\[
E(g(Y)) = E(E(X \mid Y)).
\]

\textbf{定理（重期望公式）}  
设 $(X,Y)$ 为二维随机向量，且 $E(X)$ 存在，则有
\[
E(X) = E(E(X \mid Y)).
\]

\textbf{证明（只对连续情形证明）：}

设 $(X,Y)$ 的联合概率密度为 $f(x,y)$，则有
\[
f(x,y) = f_Y(y) f_{X|Y}(x \mid y).
\]
令
\[
g(Y) = E(X \mid Y),
\]
则
\[
f_X(x) = \int_{-\infty}^{+\infty} f_Y(y) f_{X|Y}(x \mid y)\, dy
\]

于是
\[
E(X) = \int_{-\infty}^{+\infty} x f_X(x)\, dx
       = \int_{-\infty}^{+\infty} x \left[ \int_{-\infty}^{+\infty} f_Y(y) f_{X|Y}(x \mid y)\, dy \right] dx.
\]
























































\newpage

\section{随机向量的数字特征}
\begin{definition}[条件期望]
设 $(X, Y)$ 是二维离散型随机变量，对于固定的 $j$，若 $P(Y = y_j) > 0$，则定义条件概率：

\[
P(X = x_i \mid Y = y_j)
    = \frac{P(X = x_i, Y = y_j)}{P(Y = y_j)}
    = \frac{p_{ij}}{p_{\cdot j}}, \qquad i = 1, 2, \ldots
\]

若级数 $\sum g(x_i) \, p_{ij} / p_{\cdot j}$ 绝对收敛，则称

\[
\sum g(x_i)\, \frac{p_{ij}}{p_{\cdot j}}
\]

为在 $Y = y_j$ 条件下随机变量 $g(X)$ 的 \textbf{条件数学期望}。
\end{definition}
\begin{example}[条件期望的简单例子]
掷一次均匀骰子令 $X$ 为点数，定义
\[
Y =
\begin{cases}
1, & X \text{ 为偶数},\\
0, & X \text{ 为奇数}.
\end{cases}
\]
求 $E[X \mid Y = 1]$。

因为 $Y=1$ 时只有 $X = 2,4,6$，且各点概率相同，所以
\[
P(X = 2 \mid Y = 1) = P(X = 4 \mid Y = 1) = P(X = 6 \mid Y = 1) = \frac{1}{3}.
\]

由条件期望定义，
\[
E[X \mid Y = 1]
= 2 \cdot \frac13 + 4 \cdot \frac13 + 6 \cdot \frac13
= 4.
\]

因此，在已知骰子点数为偶数条件下，
\[
E[X \mid Y = 1] = 4.
\]
\end{example}
\begin{definition}
设 $(X, Y)$ 是二维连续型随机变量，联合概率密度为 $f(x,y)$，边缘密度分别为
\[
f_X(x), \qquad f_Y(y).
\]

若 $f_X(x) > 0$，则在 $X = x$ 条件下，$Y$ 的条件概率密度函数为
\[
f_{Y|X}(y \mid x) = \frac{f(x,y)}{f_X(x)}.
\]

若积分
\[
\int_{-\infty}^{\infty} g(y)\frac{f(x,y)}{f_X(x)} \, dy
\]
绝对收敛，则称
\[
\int_{-\infty}^{\infty} g(y)\frac{f(x,y)}{f_X(x)} \, dy
\]
为已知 $X = x$ 时随机变量 $g(Y)$ 的 \textbf{条件数学期望}。
\end{definition}
\begin{example}[连续型条件期望的简单例子]
设随机变量 $(X,Y)$ 在单位正方形 $[0,1]\times[0,1]$ 上均匀分布，
其联合概率密度为
\[
f(x,y) = 1, \qquad 0 \le x \le 1,~ 0 \le y \le 1.
\]

首先求边缘密度：
\[
f_X(x) = \int_0^1 f(x,y)\, dy = \int_0^1 1\, dy = 1.
\]

于是条件密度为
\[
f_{Y|X}(y \mid x) = \frac{f(x,y)}{f_X(x)} = \frac{1}{1} = 1,
\qquad 0 \le y \le 1.
\]

现在求条件期望 $E(Y \mid X = x)$：
\[
E(Y \mid X = x)
= \int_0^1 y \cdot f_{Y|X}(y \mid x)\, dy
= \int_0^1 y \cdot 1\, dy
= \frac{1}{2}.
\]

因此，不论 $X$ 取什么值，都有
\[
E(Y \mid X = x) = \frac{1}{2}.
\]
\end{example}
\begin{note}
    \[
f_{Y|X}(y \mid x)\text{ 是以 } x \text{ 为参数、以 } y \text{ 为自变量的函数}
\]

\[
E(Y \mid X = x) = \int_{-\infty}^{+\infty} y f_{Y|X}(y \mid x)\, dy \quad \text{是 } x \text{ 的函数}
\]

因为此时 \(X = x\)，故 \(E(Y \mid X = x)\) 也是 \(X\) 的函数，

\[
\text{记为 } \boxed{E(Y \mid X)}
\]
\end{note}

\begin{theorem}[全期望公式 / The Law of Total Expectation]
设 $(X,Y)$ 是二维随机变量，且 $E(g(Y))$ 存在，则有：

\begin{itemize}
    \item[(1)] （一般形式）
    \[
        E \big[\, E(g(Y)\mid X ) \,\big] = E(g(Y)).
    \]

    \item[(2)] （离散型随机变量）
    若 $X, Y$ 为离散型随机变量，且 $E(Y)$ 存在，则
    \[
        E(X) = E[E(X\mid Y)], \qquad
        E(Y) = E[E(Y\mid X)].
    \]

    \item[(3)] （连续型随机变量）
    若 $(X,Y)$ 连续且具有联合密度 $f(x,y)$，边缘密度 $f_X(x)>0$，则
    \[
        \int_{-\infty}^{+\infty} E(g(Y)\mid X=x)
        \, f_X(x) \, dx
        = E(g(Y)).
    \]
\end{itemize}

\end{theorem}

























































\newpage
\section{第三章作业}

\begin{homework}[3.1]
射击比赛中，每位运动员射 $4$ 次，每次 $1$ 发子弹。规定 4 弹都不中靶得 $0$ 分，
命中 $1$ 弹得 $10$ 分，命中 $2$ 弹得 $30$ 分，命中 $3$ 弹得 $60$ 分，命中 $4$ 弹得 $100$ 分。
若某人每次射击的命中率都是 $0.6$，求此人射击成绩（得分）的期望。
\end{homework}

\begin{homework}[3.2]
某产品的次品率为 $0.1$，检验员每天检 $4$ 次，每次任取 $10$ 件产品，
发现其中的次品数多于 $1$，就去调整设备。设备产品是否为次品相互独立，
以 $X$ 表示一天中调整设备的次数，求 $E(X)$。
\end{homework}

\begin{homework}[3.3]
将 $3$ 个球逐个独立且任意地放入编号为 $1$ 至 $4$ 的盒中，
以 $X$ 表示非空盒的最小号码，求 $E(X)$。
\end{homework}

\begin{homework}[3.4]
设离散型随机变量 $X$ 有概率分布
\[
P\!\left(X=(-1)^{i+1}\frac{3^i}{i}\right)=\frac{2}{3^i},\quad i=1,2,\ldots,
\]
说明 $X$ 的期望不存在。
\end{homework}

\begin{homework}[3.5]
设离散型随机变量 $X$ 有概率分布 $P(X=n)=2^{-n}$，$n=1,2,\ldots$，
求 $E(X)$ 和 $\mathrm{Var}(X)$。
\end{homework}

\begin{homework}[3.6]
一个有 $n$ 把钥匙的人要开他的门，他随机地独立取钥匙试开，
分别就以下两种情况求其试开次数的期望和方差：

\begin{enumerate}[(1)]
  \item 试开不成功的钥匙不从以后的选取中除去；
  \item 试开不成功的钥匙从以后的选取中除去。
\end{enumerate}
\end{homework}

\begin{homework}[3.7]
设离散型随机变量 $X$ 取正整数的概率依几何数列减少，
选择数列的首项 $a$ 及公比 $q$，使 $E(X)=10$，并计算 $P(X\le 10)$。
\end{homework}

\begin{homework}[3.8]
假设无线电台发出的呼叫信号被另一电台接收的概率为 $0.2$。
信号每隔 $5$ 秒发一次，直到收到对方的回信号为止。
若从发出信号到收到信号之间要经过 $16$ 秒，
求在双方建立起联系之前发出的呼叫信号的平均次数。
\end{homework}

\begin{homework}[3.9]
设一次试验中，事件 $A$ 出现的概率为 $p$。
以 $X$ 表示 $n$ 次重复独立试验中 $A$ 出现的次数。
离散型随机变量 $Y$ 视 $X$ 为偶数或奇数分别取 $0$ 或 $1$，求 $E(Y)$。
\end{homework}

\begin{homework}[3.10]
设每次试验中，事件 $A$ 出现的概率为 $p$，
$\overline{A}$ 出现的概率为 $q=1-p$。
若事件 $A$ 出现 $m$ 次后事件 $B$ 才出现的条件概率为
\[
G(m)=1-(1-1/\omega)^m,
\]
其中 $\omega$ 是正常数，
求为使 $B$ 出现，必须进行的独立试验次数的期望。
\end{homework}

\begin{homework}[3.11]
设连续型随机变量 $X$ 服从拉普拉斯分布，其密度为
\[
f(x)=0.5\,\exp(-|x|),\quad -\infty<x<\infty,
\]
求 $E(X)$ 和 $\mathrm{Var}(X)$。
\end{homework}

\begin{homework}[3.12]
气体分子的速度 $X$ 服从麦克斯韦分布，密度为
\[
f(x)=
\begin{cases}
\dfrac{4}{a^3\sqrt{\pi}}x^2 e^{-x^2/a^2}, & x\ge0,\\
0,& x<0,
\end{cases}
\]
求 $E(X)$ 与 $\mathrm{Var}(X)$。
\end{homework}

\begin{homework}[3.13]
设随机变量 $X$ 服从伽马分布 $\Gamma(\alpha,\beta)$，
密度为
\[
f(x)=
\begin{cases}
\dfrac{\beta^\alpha}{\Gamma(\alpha)}x^{\alpha-1}e^{-\beta x}, & x>0,\\[4pt]
0, & x\le0,
\end{cases}
\]
求 $E(X)$ 和 $\mathrm{Var}(X)$。
\end{homework}

\begin{homework}[3.14]
设随机变量 $X$ 的分布函数为
\[
F(x)=
\begin{cases}
0,& x<-1,\\
0.5+\pi^{-1}\arcsin x,& -1\le x<1,\\
1,& x\ge1,
\end{cases}
\]
求 $E(X)$ 和 $\mathrm{Var}(X)$。
\end{homework}

\begin{homework}[3.15]
设 $X$ 取非负值，证明：对任意 $\varepsilon>0$ 均有
\[
P(X<\varepsilon)\ge 1-\frac{E(X)}{\varepsilon}.
\]
\end{homework}

\begin{homework}[3.16]
设随机变量 $X$ 的期望与方差存在且取值于 $(a,b)$，证明：
\[
a\le E(X)\le b,\qquad 
\mathrm{Var}(X)\le \frac{(b-a)^2}{4}.
\]
\end{homework}

\begin{homework}[3.17]
设随机变量 $X$ 的中位数 $c$ 为满足 $F(c)=0.5$ 的数。
证明：若 $E(X)$ 存在，则
\[
E|X-c|=\min_{x\in\mathbb{R}}E|X-x|.
\]
\end{homework}

\begin{homework}[3.18]
设 $a>0$ 且 $E(e^{aX})$ 存在，证明：
\[
P(X>x)\le e^{-ax}E(e^{aX}).
\]
\end{homework}

\begin{homework}[3.19]
设二维离散型随机变量 $(X,Y)$ 的概率分布如下：
\[
\begin{array}{c|ccc}
  X\backslash Y & -1 & 0 & 1\\ \hline
  1 & 0.2 & 0.1 & 0.1\\
  2 & 0.1 & 0 & 0.1\\
  3 & 0 & 0.3 & 0.1
\end{array}
\]
求 $E(X)$，$E(Y)$ 与 $E[(X-Y)^2]$。
\end{homework}

\begin{homework}[3.20]
设二维连续型随机变量 $(X,Y)$ 的密度为
\[
f(x,y)=
\begin{cases}
12y^2,& 0\le y\le x\le1,\\[4pt]
0,&\text{其他}.
\end{cases}
\]
求 $E(X)$，$E(Y)$ 和 $E(XY)$。
\end{homework}
\begin{homework}[3.21]
设随机变量 $X$ 与 $Y$ 相互独立，分别有概率密度函数
\[
f_X(x)=
\begin{cases}
e^{-x}, & x>0,\\
0, & x\le0,
\end{cases}
\qquad
f_Y(y)=
\begin{cases}
2e^{-2y}, & y>0,\\
0, & y\le0.
\end{cases}
\]
求 $E(XY)$ 和 $E(X^2+3Y)$。
\end{homework}

\begin{homework}[3.22]
设二维连续型随机变量 $(X,Y)$ 有概率密度函数
\[
f(x,y)=
\begin{cases}
4xy\exp\{-(x^2+y^2)\}, & x>0,\,y>0,\\[4pt]
0, & \text{其他}.
\end{cases}
\]
求 $Z=\sqrt{X^2+Y^2}$ 的期望和方差。
\end{homework}

\begin{homework}[3.23]
设随机变量 $X$ 与 $Y$ 相互独立，$\mathrm{Var}(X)$ 和 $\mathrm{Var}(Y)$ 存在，证明
\[
\mathrm{Var}(XY)
=\mathrm{Var}(X)\mathrm{Var}(Y)
+[E(X)]^2\mathrm{Var}(Y)
+[E(Y)]^2\mathrm{Var}(X).
\]
\end{homework}

\begin{homework}[3.24]
设随机变量 $X\sim B(n,p)$，记 $q=1-p$，证明
\[
E(X-np)^4
=npq(p^3+q^3)+3p^2q^2(n^2-n).
\]
\end{homework}

\begin{homework}[3.25]
设二维离散型随机变量 $(X,Y)$ 有下表所示的概率分布：
\[
\begin{array}{c|ccc}
  X \backslash Y & -1 & 0 & 1 \\ \hline
  -1 & 1/8 & 1/8 & 1/8 \\
   0 & 1/8 & 0   & 1/8 \\
   1 & 1/8 & 1/8 & 1/8
\end{array}
\]
验证 $X$ 与 $Y$ 既不独立，也不相关。
\end{homework}

\begin{homework}[3.26]
设 $A$ 和 $B$ 是有正概率的事件，令
\[
X=
\begin{cases}
1, & A\text{ 出现},\\
0, & A\text{ 不出现},
\end{cases}
\qquad
Y=
\begin{cases}
1, & B\text{ 出现},\\
0, & B\text{ 不出现}.
\end{cases}
\]
若 $X$ 与 $Y$ 不相关，证明 $X$ 与 $Y$ 相互独立。
\end{homework}

\begin{homework}[3.27]
设二维离散型随机变量 $(X,Y)$ 有概率分布
\[
P(X=i,Y=j)
=\frac{r!}{i!\,j!\,(r-i-j)!}\,p_1^i p_2^j p_3^{r-i-j},
\qquad
p_1,p_2,p_3>0,\quad p_1+p_2+p_3=1,\quad i,j=0,1,2,\ldots,r,\ i+j\le r.
\]
求 $\mathrm{Cov}(X,Y)$。
\end{homework}

\begin{homework}[3.28]
在 $10$ 件产品中有 $2$ 件一级品，$7$ 件二级品，$1$ 件次品。从中不放回地抽取 $3$ 件，
用 $X$ 表示抽到的一级品数，$Y$ 表示抽到的二级品数，求 $X$ 与 $Y$ 的相关系数。
\end{homework}

\begin{homework}[3.29]
袋中有同型号小球若干，其中 $b$ 个色黑，$r$ 个色红。每次从袋中任取一球，观察其颜色后放回，
并再放入 $c$ 个同型号、同色球。令
\[
X_n=
\begin{cases}
1, & \text{第 } n \text{ 次取到红球},\\
0, & \text{第 } n \text{ 次取到黑球},
\end{cases}
\qquad n=1,2,\ldots
\]
证明：当 $m\ne n$ 时，$X_m$ 与 $X_n$ 的相关系数为 $c/(b+r+c)$。
\end{homework}

\begin{homework}[3.30]
设二维连续型随机变量 $(X,Y)$ 有概率密度函数
\[
f(x,y)=
\begin{cases}
1, & |y|<x,\ 0<x<1,\\[4pt]
0, & \text{其他}.
\end{cases}
\]
求 $\mathrm{Cov}(X,Y)$。
\end{homework}

\begin{homework}[3.31]
设二维连续型随机变量 $(X,Y)$ 有概率密度函数
\[
f(x,y)=
\begin{cases}
\dfrac{1}{8}(x+y), & 0\le x\le2,\ 0\le y\le2,\\[4pt]
0, & \text{其他}.
\end{cases}
\]
求 $\mathrm{Var}(X+Y)$、$\mathrm{Cov}(X,Y)$ 及 $\rho_{X,Y}$。
\end{homework}

\begin{homework}[3.32]
设 $X,Y,Z$ 是随机变量，且
\[
E(X)=E(Y)=1,\quad E(Z)=-1,\quad 
\mathrm{Var}(X)=\mathrm{Var}(Y)=\mathrm{Var}(Z)=1,
\]
\[
\rho_{X,Y}=0,\quad \rho_{X,Z}=0.5,\quad \rho_{Y,Z}=-0.5.
\]
求 $E(X+Y+Z)$ 和 $\mathrm{Var}(X+Y+Z)$。
\end{homework}

\begin{homework}[3.33]
设二维连续型随机变量 $(X,Y)$ 有概率密度函数
\[
f(x,y)=
\begin{cases}
0.5\sin(x+y), & 0\le x\le0.5\pi,\ 0\le y\le0.5\pi,\\[4pt]
0, & \text{其他}.
\end{cases}
\]
求 $\mathrm{Cov}(X,Y)$ 与 $\rho_{X,Y}$。
\end{homework}

\begin{homework}[3.34]
设 $n>m$，随机变量 $X_1,X_2,\ldots,X_{n+m}$ 相互独立，且具有相同的分布函数，
$\mathrm{Var}(X_1)$ 存在。在
\[
Y=\sum_{i=1}^{n} X_i,\qquad Z=\sum_{k=1}^{m} X_{n+k}
\]
下，求 $\rho_{Y,Z}$。
\end{homework}

\begin{homework}[3.35]
在长为 $l$ 的线段上任取两点，求两点间距离的期望与方差。
\end{homework}

\begin{homework}[3.36]
设二维连续型随机变量 $(X,Y)$ 服从非退化的二维正态分布
\[
N(0,1;0,1;r),\quad |r|<1.
\]
证明：
\[
E[\max(X,Y)] = \sqrt{\frac{1-r}{\pi}}.
\]
\end{homework}

\begin{homework}[3.37]
设每天进入货站的货物件数 $N$ 有如下概率分布：
\[
\begin{array}{c|cccccc}
N=n & 10 & 11 & 12 & 13 & 14 & 15 \\ \hline
P(N=n) & 0.05 & 0.10 & 0.10 & 0.20 & 0.35 & 0.20
\end{array}
\]
并设每天到达的货物中次品的概率均为 $0.10$。
以 $X$ 表示每天进货中次品的件数，求 $E(X)$。
\end{homework}

\begin{homework}[3.38]
设随机变量 $X_1 \sim P(\lambda_1)$，$X_2 \sim P(\lambda_2)$，且二者独立。

\begin{enumerate}[(1)]
  \item 证明：当 $n \ge 2$ 时，有
  \[
  P(X_1=i \mid X_1+X_2=n)
  = C_n^i
    \left(\frac{\lambda_1}{\lambda_1+\lambda_2}\right)^i
    \left(\frac{\lambda_2}{\lambda_1+\lambda_2}\right)^{n-i},
  \qquad i=0,1,2,\ldots,n.
  \]
  \item 求 $E(X_1\mid X_1+X_2=n)$ 和 $E(X_2\mid X_1+X_2=n)$。
\end{enumerate}
\end{homework}

\begin{homework}[3.39]
设二维连续型随机变量 $(X,Y)$ 有概率密度函数
\[
f(x,y)=
\begin{cases}
8xy, & 0<y<x<1,\\[4pt]
0, & \text{其他}.
\end{cases}
\]
求条件期望 $E(X\mid Y=y)$ 与 $E(Y\mid X=x)$。
\end{homework}

\begin{homework}[3.40]
已知随机变量 $Y\sim \Gamma(\alpha,\beta)$，且在 $Y=y$ 条件下，
$X$ 的条件概率密度函数为
\[
f_{X\mid Y}(x\mid y)=
\begin{cases}
y\exp(-yx), & x>0,\\[4pt]
0, & \text{其他}.
\end{cases}
\]
\begin{enumerate}[(1)]
  \item 求在 $X=x$ 条件下，$Y$ 的条件概率密度函数 $f_{Y\mid X}(y\mid x)$；
  \item 求 $E(Y\mid X=x)$。
\end{enumerate}
\end{homework}